\section{Introduction}

So welcome everybody to the course on the geometric anatomy of theoretical physics. You, being advanced students, know that theoretical physics is all about casting our concepts about the real world into rigorous mathematical form, for better or worse. But theoretical physics doesn't do that for its own sake; it does so in order to fully explore the implications of what our concepts about the real world are. So to a certain extent, the spirit of theoretical physics can be cast into the words of Wittgenstein, who said: ``What we cannot speak about clearly, we must pass over in silence.''

Well, so indeed, if we have concepts about the real world and it's not possible to cast them into precise mathematical form, that usually is an indicator that some aspects of these concepts have not been well understood. So theoretical physics aims at casting these concepts into mathematical language. But then mathematics is just that---it's just a language. And if we want to extract physical conclusions from this formulation, we must interpret the language, and that's not the purpose or the task of mathematics; that's the task of physicists. And that is where it gets difficult again. But again, mathematics is just a language, and to again use Wittgenstein, he says: ``The theorems of mathematics all say the same, namely nothing.'' Okay, and so what does he mean by that? Well, obviously he doesn't mean that mathematics is useless. He refers to the fact that if we have a theorem of the type $A$ if and only if $B$, $A$ and $B$ being propositions, then obviously $B$ says nothing else than $A$ does, and $A$ says nothing else than $B$ does---it's a tautology. Okay, now of course, mathematically, logically, it's a tautology, but psychologically, or for understanding of $A$, it may be very useful to have a reformulation of $A$ in terms of $B$, right?

So with this understanding that mathematics just gives us a language for what we want to do, the idea of this course is to provide proper language for theoretical physics, and certainly for the courses you have heard so far. So the aims of the course: the overarching aim is to provide proper mathematical language for classical mechanics, electromagnetism, quantum mechanics, and statistical physics. Now, these subjects you studied in the modules TP1 to TP4 here at FAU, or at many other universities. And obviously we're not going to revise all the mathematics that are needed for these four subjects, but rather we will develop mathematics from a higher point of view that we can now approach, given that you have some prior knowledge of these subjects.

So far you have certainly used a fair amount of analysis in order to study these subjects, further a fair amount of algebra---mainly linear algebra, but also nonlinear algebra at times---and geometry in one or the other form. So very often in classical mechanics you would appeal to geometric intuition in setting up a problem or finding the solution to a problem. The same goes for electromagnetism. Algebra you mainly used in quantum mechanics; that's the modern way to do quantum mechanics, by means of linear algebra. And in statistical physics you use some geometry, maybe without really realizing it, by looking at convex cones of states and so on.

Now if you take these four fields and you try to be pedagogical, if you draw circles around them: classical mechanics and electrodynamics certainly use a fair amount of analysis and geometry, and so does general relativity, for those of you who took that course or are about to take it. Quantum mechanics lies certainly well in the intersection of analysis and algebra, usually called functional analysis. And algebra and geometry---let's do statistical physics justice---it more or less lies there. Of course this is not a pure picture; this is mainly where the intersections of these fields come to play. And here in the middle, where these three fields intersect, we have differential geometry as a mathematical subject. And it's differential geometry that this course is mainly about. So the mathematics of this course will be that part of the mathematics of these different theoretical subjects where these three fields intersect. That's the rough idea of this course.

Now the structure of this course: indeed, the development of differential geometry that we're about to start begins with sets. We'll talk about set theory, because any space---in classical mechanics the space of configurations or the phase space, in electromagnetism the space in which you work, or in quantum mechanics the physical space over which you construct your Hilbert space, and so on---all these spaces are certainly sets of points. That's the coarsest structure of such a space: it is a set. But what precisely is a set? Well, one can delude oneself into thinking one understands a set by saying a set is a collection of elements, but that just raises the question what a collection is and what elements are. So we certainly need to do better.

And there is a fundamental problem: if we start writing a book about mathematics whose pages are all empty yet, what could the first definition be? For a definition, you need notions that you already have in order to define a new notion. But if you don't have any notion yet, how do you start? Well, the trick is to start axiomatically. And so we'll have to write down axiomatic set theory. But that then raises the question: in what language would you possibly do that? So actually, before we come to set theory, we need another building block down here, and that will be logic. We will deal with propositional and predicate logic first; that will define our language, and then we'll be able to write down the axioms of set theory.

Now once we have the axioms of set theory and we have our beautiful sets, the question is: what can we do with them? And one thing we certainly need, at least for these four subjects, is a notion of continuity. For instance, in classical mechanics, famously, particles don't jump---the particle trajectory doesn't stop somewhere and start somewhere else again; it's continuous. So you need a notion of continuity, and such a notion you get by distinguishing certain subsets of the basic set you have, and that leads to the notion of topology. So we'll deal with topological spaces and all the notions associated with that.

Now there is an abundance of topological spaces---such an abundance that you cannot even classify them---and so we will need to restrict somewhat to interesting topological spaces. Well, topologists may disagree that only those I'm going to write down are the interesting ones, but from the point of view of theoretical physics, immediately interesting are the topological manifolds. Topological manifolds, which we're going to define as topological spaces that locally around a point look like some $\mathbb{R}^d$.

Now that's not enough. If you have topological manifolds, you also need in physics, especially in classical physics, a notion of derivative. You need to be able to talk about the velocity of a curve, and that cannot be done in topological manifolds. What you need there are differentiable manifolds, and that's an entirely new beast. We will need to introduce the technology of charts, and the technology of charts will for the first time make contact to much of the formalism you already used in the standard courses. In fact, all the standard courses have been formulated in charts so far.

I'll explain later: a chart is like a sailor uses a chart of the world. You travel across the globe, and if you want to know where you are and where you want to sail, you take out an atlas. In the atlas you'll turn the pages and you'll find the appropriate chart around the point where you are in the real world. You can navigate in the chart, and the chart should be a one-to-one picture of the world locally, because if you then move in the world such that in the chart you would leave the page, you better turn the page. And you'll find the other page that slightly overlaps with the chart in the atlas, so that you can then navigate on with the other chart. And very much the same thing---looking for another chart that overlaps and understanding how the overlap works---that's the whole philosophy behind manifolds, and certainly behind differentiable manifolds. Because differentiation: the charts are parts of $\mathbb{R}^d$, like $\mathbb{R}^2$ for instance, and so you understand differentiation in the charts. And we will lift the notion of differentiation that we understand well in the charts; we'll lift it to the general concept of differentiation on the manifold. That's already the entire philosophy.

So we will certainly come to that. Now once we have differentiable manifolds, maybe already before, we'll come to a notion which, well, in some sense is yet another structure, but it's bundles. So the idea of bundles is you take a manifold and you take another manifold, and to each point of a given manifold you attach the other manifold. You build a bundle: you could have a vector bundle; at each point you attach the tangent space, and so on. And bundles are a very important concept in geometry, in differential geometry, and indeed much of modern physics can be understood---or can be cast into---the language of bundles. So this is true for particle physics, but also for gravity and so on; certainly for quantum mechanics.

So actually, about here, already here, we will learn that the notion of a position vector should be forbidden. There is no such thing as a position vector; doesn't matter how often it's repeated, there is no such notion. Now, with bundles you learn there is no such thing as a wave function in quantum mechanics. You get away with the notion of position vector unless you look too closely, and then it fails. And the same is true in quantum mechanics: you get away with the notion of wave function, but then it fails. In fact, a position vector is just a collection of coordinates, but they do not constitute a vector. And the wave function in quantum mechanics is not a function at all; it's a section in a complex line bundle over the physical space. You cannot understand this now, but you will at the end of the course. So these notions are needed if you really want to say what's going on.

Now, okay, so once we have bundles, or maybe also as an alternative to bundles sometimes---so don't take the picture too seriously, it's a rough idea of how this building is constructed---we will have geometry. Geometry comes in the form of so-called tensor fields on this structure. So that could be symplectic geometry; that would be the geometry of a classical phase space in mechanics. But actually it's a much bigger subject, because for virtually all dynamical systems of interest to us, the phase space is equipped with symplectic geometry. There may be a metric geometry; this is certainly the case in electromagnetism. Due to a certain extent, it's needed in geometry just in the volume form. And certainly it's the case in relativity---special and general relativity---you need a so-called Lorentzian metric. This is a structure that sits on top here, but you need to understand more or less all of this in order to really professionally deal with geometry.

Well, and then to complete this: I promised you that all these different subjects you learned about before---so up to here it's just physical mathematics, but then up here it's physics. We have classical mechanics, electromagnetism, quantum mechanics, statistical physics, and special relativity and general relativity. All of these physics can be understood, and can be properly understood, by mastering all this mathematics. And again, the purpose of this course is to develop this mathematics and then to show, mainly by examples, how the insights you gain from here play out at the level of physics. That's the idea of this course.
